\section{Bar--cobar construction}
\label{ch:bar-cobar}

This chapter compares the cubic action with the bar--cobar construction of Chiaffrino, Mamade and Zwiebach, which rewrites the equations of motion of any $\Loo$ field theory as a quadratic equation on an enlarged space of multi-string fields but gives no action for it.
The paper keeps a shortened comparison in three subsections: their construction in our notation, the canonical solutions and the gauge symmetry, and the map to the cubic action.
This chapter reproduces the comparison section of the paper as it stood before the paper was shortened, in two blocks: the first holds the construction and the identification of the fields on the canonical sector; the second holds the subsection on the cubic action as a truncation of their equations, which is kept here in full, the gauge symmetry, the additional structure of the cubic action and the limits of the comparison.
The identification of their canonical equations with the multiplier form, which stood between the two blocks, is reproduced in \cref{sec:mult:barcobar} and is not repeated.
The last section, on the two obstructions to an action and on their scalar example, is written for this document from the derivation and gives the degree count, the self-adjointness test and the scalar computation that the paper states in a few sentences.

The comparison establishes the following statements, with the status of each.
All the finite-dimensional checks are made in models with commuting fields, where no Koszul sign arises and no type-II field is realized, so that the gauge statements are blind to signs; the solutions for $\Psi$ of the two sets of equations coincide order by order in $\gs$ and modulo the kernel of the quartic form, the auxiliary fields agreeing with the canonical values only at leading order.

\begin{itemize}
  \setlength{\itemsep}{1ex}

  \item \emph{The comparison as a whole.}

    \status{mixed}\record{024-hs-vs-bar-cobar (R34)}

  \item \emph{Fields and canonical equations.} The multiplier form is not the cotangent action of their type-I sector as such, since the split $(2, 2)$ at length four and the components of $B \Phi_1 = 0$ with two or more output strands have no Euler--Lagrange counterpart; it is the cotangent action of the chain alone, with $\Psi$ not doubled; the multipliers $\Theta_k$ have ghost number $3 k$ for the pairing $\pair[k]{\cdot}{\cdot}$.

    \status{Proven} for the term-by-term comparison and the ghost number; \status{Numerical} for the cotangent action of the chain, in a model with commuting fields.

  \item \emph{The cubic action as a truncation.} At quartic order, the identity~\eqref{p:eq:barcobar:identity} between the chain action with its multiplier on-shell and the quartic term of the cubic action holds off-shell.

    \status{Proven}, and checked in a model with commuting fields.

    The completion at order $\gs^3$ is the quintic term of the cubic action, and the action with $\Sigma_2$ eliminated differs from the original one by a term of order exactly $\gs^4$.

    \status{Conditional(PH1, PH5) -- Numerical}, in a model with commuting fields and a generic invertible quartic form.

    The cubic action is not a consistent truncation of their equations.

    \status{Argued}, by the structural argument of the block.

  \item \emph{Gauge symmetry.} Restricted to the canonical sector, their gauge transformation of $\Sigma_2$ is the composite law plus a term proportional to the equation of motion of $\Psi$, \eqref{p:eq:barcobar:gauge-sigma}.

    \status{Numerical}, in a model with commuting fields.

    Their independent two-string parameter $\hat\Lambda_{1, 2}$ is not a symmetry of the cubic action on type-I fields.

    \status{Argued}, from a ghost-number saturation argument on the gauge laws of the auxiliary field; the model check corroborates it only weakly.

    The parameters that preserve the canonical form are the single-string ones, \eqref{p:eq:barcobar:canonical-parameters}.

    \status{Argued}: their theorem, transcribed to the symmetric coalgebra and not re-derived.

  \item \emph{The two obstructions.} The degree count~\eqref{eq:bc:pairing-degree} of the length-diagonal pairings, and the realization of the second obstruction in sharpened form on a model, in \cref{sec:bc:obstructions}.

    \status{Proven} for the degree count; \status{Numerical} for the second obstruction on the model.

  \item \emph{The scalar example} of $\phi^3 + \phi^4$ theory, reproduced by the multiplier form with three multipliers, in \cref{sec:bc:scalar}.

    \status{Numerical}, on a four-site lattice.
\end{itemize}

The first block is the introduction of the paper's comparison section, the recall of their construction in our notation (factor number, length, the degree~\eqref{p:eq:barcobar:degree} of a type-I field in terms of its ghost number, the canonical equations and their unique solution), the two obstructions as they state them, and the identification $\Sigma_r = r! \, \Phi_{1, r}$ of their type-I fields with the auxiliary fields.
It uses the notation of the paper throughout: $\Hn{r}$ for the $r$-string space, $\pair[r]{\cdot}{\cdot}$ for the $r$-string pairing, $\vertex{n}$ for the vertices and $\Theta_k$ for the multipliers.

\subsection{Bar–cobar construction}
\label{p:sec:bar-cobar}


Chiaffrino, Mamade and Zwiebach have shown that the equations of motion of any $\Loo$ field theory, closed string field theory included, can be rewritten as a quadratic equation on an enlarged space of multi-string fields~\cite{Chiaffrino:2026:UniversalQuadraticField}.
Their proofs are written for $A_\infty$ algebras. For $\Loo$ algebras, the equivalence follows from a general theorem on quasi-isomorphic algebras.
They give no action for that equation and leave its existence open.
The enlarged space contains one multi-string field for every number of strands, which they call type I, besides fields valued in formal products of such multi-string spaces, their type II (both are defined below).
Its $2$-strand component is the space in which $\Sigma_2$ lives, up to a quotient, so that we can compare the two constructions field by field.
In this section, we recall their construction in our notation, we identify the fields and the equations which the two constructions share, and we state in what sense neither construction is a truncation of the other.
The comparison is classical and algebraic.


\subsubsection{The bar–cobar construction}
\label{p:sec:barcobar:construction}

The products $\ell_n$ of closed string field theory form an $\Loo$ algebra~\cite{Zwiebach:1993:ClosedStringField,Erbin:2021:StringFieldTheory}.
The bar construction packages all of them into a single coderivation $B$ of the symmetric tensor coalgebra built on the state space, and the cobar construction then builds the free Lie algebra on this coalgebra.
The result is a differential graded Lie algebra, quasi-isomorphic to the original $\Loo$ algebra~\cite{Chiaffrino:2026:UniversalQuadraticField}.
Its Maurer--Cartan equation is quadratic because a differential graded Lie algebra has no product beyond the bracket. The price is the size of the field space.

The field $\Phi$ of the enlarged theory has components labeled by a \emph{factor number} $f$ and by a \emph{length}, the number of strands.
The components of factor number one, which they call type I, are one multi-string field per length:
\begin{equation}
  \label{p:eq:barcobar:type-I}
  \Phi
    = \Phi_1 + \Phi_2 + \cdots,
  \qquad
  \Phi_1
    = \sum_{r \ge 1} \Phi_{1, r},
  \qquad
  \Phi_{1, r}
    \in \Hn{r},
  \qquad
  \Ngh(\Phi_{1, r})
    = 2 r,
\end{equation}
and the length-one component $\Phi_{1,1}$ is the string field itself.
We describe their construction in our notation, translating their homological degree into ghost number: a type-I field of length $r$ and total ghost number $N$ has degree
\begin{equation}
  \label{p:eq:barcobar:degree}
  \deg \Phi_{1, r}
    = 2 r - 1 - N
\end{equation}
in their enlarged space.
The components of factor number $f \ge 2$, their type II, are independent fields valued in the formal products $\Hn{r_1} \otimes' \cdots \otimes' \Hn{r_f}$ generated by the free Lie algebra of the cobar step.
A type-II component with $f$ factors and total length $r$ has ghost number $2 r + 1 - f$, which differs from the ghost number $2 r$ of the auxiliary field with the same number of strands.
The equation of motion reads:
\begin{equation}
  \label{p:eq:barcobar:mc}
  (B + \Delta_c) \Phi + \Phi \otimes' \Phi
    = 0,
\end{equation}
where $\Phi \otimes' \Phi = \tfrac{1}{2} [\Phi, \Phi]$ for a field of their degree $-1$.
All the data of the theory sit in $B$: on a type-I field of length $r$ it applies the products $\gs^{n-1} \ell_n$, for every $n \ge 1$, to $n$ of the strands and sums over the choice of the strands.
It lowers the length by $n - 1$ and preserves the factor number.
The two other operations are universal and raise the factor number by one.
The coproduct $\Delta_c$ splits a type-I field of length $n$ into the pairs of lengths $(i, n - i)$, by summing over the ways of distributing the $n$ strands, and $\otimes'$ juxtaposes two fields.

Since $B$ preserves the factor number and the two other operations raise it by one, an ansatz with type-I components only, $\Phi = \Phi_1$, reduces~\eqref{p:eq:barcobar:mc} to two equations:
\begin{equation}
  \label{p:eq:barcobar:canonical-eqs}
  B \Phi_1
    = 0,
  \qquad
  \Delta_c \Phi_1 + \Phi_1 \otimes' \Phi_1
    = 0,
\end{equation}
of factor numbers one and two, all higher factor numbers vanishing identically.
The second equation is a factorization condition: component by component, for every length $n$ and every split $1 \le i \le n - 1$:
\begin{equation}
  \label{p:eq:barcobar:factorization}
  \binom{n}{i} \, \Phi_{1, n}
    = \Phi_{1, i} \otimes \Phi_{1, n - i},
\end{equation}
as tensors with $n$ strands, the right-hand side being required to be symmetric under all permutations of its strands.
Its unique solution is the exponential of the string field:
\begin{equation}
  \label{p:eq:barcobar:canonical}
  \Phi_{1, r}
    = \frac{\Psi^r}{r!},
\end{equation}
which satisfies~\eqref{p:eq:barcobar:factorization} for every split.
On this solution, the first equation of~\eqref{p:eq:barcobar:canonical-eqs} reduces to the equation of motion $\mathcal E_\Psi = 0$ of the original theory, with $\mathcal E_\Psi$ given by~\eqref{p:eq:hs1:eom-psi}.
Indeed, the components of $B \Phi_1$ with $L \ge 2$ output strands are $\mathcal E_\Psi$ tensored with $\Psi^{L-1} / (L-1)!$.

They call these solutions \emph{canonical}.
They prove, order by order in $B$, that every solution of~\eqref{p:eq:barcobar:mc} is gauge equivalent to a canonical one~\cite{Chiaffrino:2026:UniversalQuadraticField}.
The gauge classes of solutions of the quadratic equation are therefore in one-to-one correspondence with those of the original theory.
They also show that, among the parameters of factor number one, those which preserve the canonical form generate the gauge transformations of the original theory.
In the $\Loo$ case, where the bar construction uses the symmetric coalgebra, the same argument gives a single-string parameter $\Lambda \in \Hs$ of ghost number one.

Their construction gives equations but no action, and they identify two obstructions~\cite{Chiaffrino:2026:UniversalQuadraticField}.
First, an action for a cyclic $\Loo$ algebra needs a non-degenerate bilinear form of a definite degree, and they argue that the enlarged space admits none.
In their scalar field example every vector has negative degree, so that nothing is available to pair against.
Second, they argue that a form diagonal in the factor number and the length is unlikely to be cyclic, since $\Delta_c$ and $\otimes'$ raise the factor number.
Indeed $\langle \Phi_2, \Phi_1 \otimes' \Phi_1 \rangle$ can be nonzero while $\langle \Phi_1, \Phi_2 \otimes' \Phi_1 \rangle$ vanishes, and for the same reason $\Delta_c$ is unlikely to be self-adjoint.
They leave the question open.

We can easily write an action whose equations contain theirs by doubling every field with a multiplier. The question they leave open is the existence of an action on their field space itself.
They also note that simpler constructions should be possible with more knowledge of the theory, and give the Hubbard--Stratonovich transformation of a quartic interaction with one bilocal field as an example.
For closed string field theory this example is the earlier construction of~\cite{Erbin:2018:CubicClosedString}, which the paper extends to all orders.

\subsubsection{Identification on the canonical sector}
\label{p:sec:barcobar:canonical}

The type-I component of length $r$ lives in $\Hn{r}$ with ghost number $2 r$, which is the space and the ghost number of the auxiliary field $\Sigma_r$.
The identification is simplest when we absorb the factorials:
\begin{equation}
  \label{p:eq:barcobar:identification}
  \Sigma_r
    := r! \, \Phi_{1, r},
  \qquad
  r \ge 1,
\end{equation}
so that $\Sigma_1 = \Phi_{1, 1} = \Psi$, the canonical solution is $\Sigma_r = \Psi^r$, and the factorization conditions~\eqref{p:eq:barcobar:factorization} read $\Sigma_n = \Sigma_i \otimes \Sigma_{n-i}$ for every split.


The second block follows, in the paper, the identification of the multiplier form, whose action~\eqref{p:eq:barcobar:multiplier-action} it uses from its first paragraph.
It relates the cubic action to that action at quartic order, lists what the truncation keeps and drops, compares the gauge symmetries, and states in a few sentences how the cubic action evades the two obstructions; the next section gives these sentences in full.


\subsubsection{The Hubbard--Stratonovich action as a truncation}
\label{p:sec:barcobar:truncation}

The cubic action of \cref{p:sec:hs} keeps a finite number of auxiliary fields at each order, $\Sigma_k$ for $2 k \le n$ at order $\gs^{n-2}$, and carries no multipliers.
We relate it to~\eqref{p:eq:barcobar:multiplier-action} at quartic order, where the relation is an identity.
We keep $\Psi$, $\Sigma_2$ and $\Theta_2$, and we replace $\Sigma_3$ in the cubic vertex by its value $\Psi^3$ on the diagonal, so that $\Theta_3$ never appears.
We also replace $\Sigma_4$ in the quartic vertex by $\Sigma_2 \otimes \Sigma_2$, the split $(2, 2)$ of~\eqref{p:eq:barcobar:factorization}, so that this vertex becomes quadratic in $\Sigma_2$:
\begin{equation}
  \label{p:eq:barcobar:chain-action}
  \caction[\mathrm{chain}]
    := S_2 + S_3
      + \frac{\gs^2}{4!} \, \vertex{4}(\Sigma_2^2)
      + \pair[2]{\Theta_2}{\Sigma_2 - \Psi^2}.
\end{equation}
The multiplier form does not impose this split, so that~\eqref{p:eq:barcobar:chain-action} is not simply~\eqref{p:eq:barcobar:multiplier-action} truncated to $k \le 2$.
The $\Sigma_2$ equation of~\eqref{p:eq:barcobar:chain-action} fixes the multiplier:
\begin{equation}
  \label{p:eq:barcobar:composite-multiplier}
  \pair[2]{\Theta_2}{\cdot}
    = - \frac{\gs^2}{12} \, \vertex{4}(\Sigma_2, \cdot).
\end{equation}
Substituting this value back into the action, off-shell in $(\Psi, \Sigma_2)$, we get:
\begin{equation}
  \label{p:eq:barcobar:identity}
  \frac{\gs^2}{4!} \, \vertex{4}(\Sigma_2^2)
    - \frac{\gs^2}{12} \, \vertex{4}(\Sigma_2, \Sigma_2 - \Psi^2)
    = - \frac{\gs^2}{4!} \, \vertex{4}(\Sigma_2^2)
      + \frac{2 \gs^2}{4!} \, \vertex{4}(\Sigma_2, \Psi^2)
    = \caction[4],
\end{equation}
the quartic term of the cubic action of \cref{p:sec:hs:quartic}, with its ratio $2 : 1$ between the source and the quadratic term.
The ratio follows from the structure alone.
For any term $F(\Sigma_2)$ quadratic in $\Sigma_2$, the multiplier fixed by the $\Sigma_2$ equation turns $F(\Sigma_2) + \pair[2]{\Theta_2}{\Sigma_2 - \Psi^2}$ into the completed square $F(\Psi^2) - F(\Sigma_2 - \Psi^2)$.
For $F = \tfrac{\gs^2}{4!} \, \vertex{4}$ this is $S_4 + \Delta S_4$ of~\eqref{p:eq:hs4:S4tilde}.
The chain action~\eqref{p:eq:barcobar:chain-action} and $S_2 + S_3 + \caction[4]$ have the same solutions for $\Psi$, and for $\Sigma_2$ modulo $\ker \sprod{2}{2}$.
They are not related by a field redefinition, since $\Theta_2$ is an independent field in the first and absent from the second.

The identity~\eqref{p:eq:barcobar:identity} shows where the cubic action finds the pairing that their space lacks.
The multiplier form adds a dual tower. The Hubbard--Stratonovich action instead uses the quartic vertex itself, a Grassmann-even symmetric form on the $2$-string space supplied by the theory.
This form needs no odd symplectic structure, because $\Sigma_2$ has no $Q_B$ kinetic term, and it is degenerate exactly on $\ker \sprod{2}{2}$, which is why $\Sigma_2$ lives on the quotient of \cref{p:sec:hs:quartic}.

Their scalar field example shows the difference with the string case.
There the quartic vertex is the contact term $\int \phi^4$, which vanishes on every off-diagonal component of a bilocal field $\Sigma_2(x, y)$.
The quotient therefore reduces $\Sigma_2$ to the local field $\Sigma_2(x, x)$ of the usual Hubbard--Stratonovich transformation.
In string field theory the quartic vertex integrates over a region of moduli space, its kernel is the one of \cref{p:sec:hs:quartic}, and $\Sigma_2$ remains a genuine $2$-string field.

At quintic order the only rewriting of $\vertex{5}(\Sigma_5)$ which is cubic in $\Psi$ and $\Sigma_2$ is $\vertex{5}(\Sigma_2^2, \Psi)$, with the coefficient $\gs^3 / 5!$ of \cref{p:sec:hs:quintic}.
The price is visible there.
From the quintic order on, we assume, as in \cref{p:sec:hs:quartic,p:sec:hs:quintic}, that the quadratic form $\vertex{4}$ is invertible on the quotient by its kernel, level by level, and that the quintic vertex vanishes on that kernel.
The kinetic operator of $\Sigma_2$ is replaced by the dressed operator $\kin{2}{\Psi}{3}$, the solution $\Sigma_2^{(3)*} = \Psi^2 + O(\gs)$ leaves the diagonal, and the eliminated action reproduces $S_{\le 5}$ up to a term of order $\gs^4$ which the next order absorbs.
In their construction the conditions $\Sigma_k = \Psi^k$ are exact equations at every order and $B$ acts by the elementary products on the strands.
The truncation trades this exactness for a cubic action in a finite number of fields, and re-absorbs the mismatch order by order through the reduced vertices $\dvertex{n}$ of \cref{p:sec:all-orders}, under the assumptions stated there.

We can now list the parts of their structure that are kept and those that are dropped.
The type-II fields are dropped altogether, which costs nothing for canonical solutions by their theorem.
The type-I fields of length above $n / 2$ at order $\gs^{n-2}$ are replaced by tensor products of lower-order fields, $\Psi$ included, given by the factorization conditions~\eqref{p:eq:barcobar:factorization}, rather than set to zero, so that every vertex becomes quadratic or cubic in the fields kept.
The multiplier tower is replaced by the quadratic forms of the $\Sigma_k$, built from $\vertex{2k}$ and reduced by the lower-order fields and exchanges (\cref{p:sec:all-orders}).
The factorization condition for $\Sigma_2$ becomes the equation of motion $\sprod{2}{2}(\Sigma_2 - \Psi^2) = O(\gs)$, so $\Sigma_2 = \Psi^2$ holds only modulo $\ker \sprod{2}{2}$ and at leading order.
The components of $B \Phi_1 = 0$ with several output strands are not imposed.

In this precise sense, the cubic action is not a consistent truncation of their equations.
Fields are substituted rather than set to zero, the pairing is changed, and the solution for $\Sigma_2$ leaves the diagonal at order $\gs$, while the original action is recovered order by order.
Conversely their theory is not a truncation of ours, which has an action, no type-II fields and a finite number of auxiliary fields at each order.
The cubic action has the field content of their canonical type-I sector up to the largest length kept at each order.
Its equations of motion have the same solutions for $\Psi$ as theirs in the canonical gauge $\Phi_{f \ge 2} = 0$, order by order in $\gs$ and modulo the kernels of the quadratic form and of the nested operators of the auxiliaries, under the assumptions of \cref{p:sec:all-orders}. The auxiliary fields agree with $\Sigma_k = \Psi^k$ only at leading order and modulo these kernels.
Indeed, our equations of motion reproduce the original one in this sense, their canonical equations reduce to it as shown above, and their theorem is needed only to reach the canonical gauge.

\subsubsection{Gauge symmetry}
\label{p:sec:barcobar:gauge}

Their gauge transformation,
\begin{equation}
  \label{p:eq:barcobar:gauge}
  \delta \Phi
    = (B + \Delta_c) \hat\Lambda + [\Phi, \hat\Lambda],
\end{equation}
has one parameter $\hat\Lambda_{f, r}$ for every factor number $f$ and every length $r$.
In particular it contains an independent $2$-string parameter $\hat\Lambda_{1, 2} \in \Hn{2}$ of ghost number three, of the same degree as the parameter $\Lambda_2$ of \cref{p:sec:gauge}.
  Within the classes examined there, every theory of which the cubic action is a gauge-fixed form is a trivial enlargement; we have not worked out how $\hat\Lambda_{1, 2}$ relates to these classes.
Restricted to the canonical sector and to parameters of factor number one, the transformation preserves the canonical form only for
\begin{equation}
  \label{p:eq:barcobar:canonical-parameters}
  \hat\Lambda_{1, r}
    = \frac{1}{(r - 1)!} \, \Lambda \otimes \Psi^{r - 1},
\end{equation}
symmetrized over the strands, with a single-string $\Lambda$, by the symmetric form of their result quoted above.
Within the canonical sector, the $2$-string parameter is therefore fixed by the single-string one.
An independent $\hat\Lambda_{1, 2}$ generates type-II components and leaves the canonical sector.
On $\Psi$ the transformation is the gauge transformation of the original theory, and on $\Sigma_2$ it reads:
\begin{equation}
  \label{p:eq:barcobar:gauge-sigma}
  \delta_\Lambda \Sigma_2
    = 2 \, \delta_\Lambda \Psi \otimes \Psi
      + 2 \, \mathcal E_\Psi \otimes \Lambda,
\end{equation}
where $\mathcal E_\Psi$ is the equation of motion of the original theory and the products are understood inside a symmetric vertex as in \cref{p:sec:hs:quartic}.
We have fixed the relative sign of the second term only in finite-dimensional models with commuting fields, of the kind described in \cref{p:app:finite-models}, where no Grassmann sign arises.
The first term is the composite law of \cref{p:sec:hs:quartic}.
The second term vanishes on-shell, and it has to be there: their transformation is a symmetry of their equations, which have no action, and preserves the diagonal $\Sigma_2 = \Psi^2$ on-shell only, whereas the composite law preserves it off-shell.
  The two laws differ by $2 \, \mathcal E_\Psi \otimes \Lambda$, a term proportional to the equation of motion of the original theory, so that they agree on-shell.

The independent parameter $\hat\Lambda_{1, 2}$ exists in their theory, where it is used up in reaching the canonical gauge.
By their theorem, under the same assumptions, the type-II components of any solution can be set to zero by a gauge transformation, after which only the single-string parameters remain.
The gauge symmetry of the cubic action is thus the residue of theirs in the canonical gauge.

\subsubsection{Additional structure of the cubic action}
\label{p:sec:barcobar:structure}

Their quadratic equation is universal. The cubic action, instead, is built from the vertices of the theory, needs its quartic form to be invertible on the quotient by its kernel, and adds two things.
First, it is an action. Its equations of motion are quadratic in the fields, like theirs, but they derive from a variational principle.
It comes with the Feynman rules in which the exchange of $\Sigma_2$ reproduces the quartic vertex, with a Gaussian integral over the auxiliary fields, with a Batalin--Vilkovisky master action and a spectrum (\cref{p:sec:gauge}), and with classical solutions.
Second, it is gauge invariant to the order of the truncation, with a transformation of the auxiliary fields induced by that of $\Psi$.
This is the composite law at quartic order, corrected from the quintic order on by a rotation of the fluctuation (\cref{p:sec:gauge}).
Closing the symmetry at each order requires the reduced vertices, with a finite number of auxiliary fields at each order rather than their whole tower.

We now return to their two obstructions.
In our grading, their counting argument does not apply as such, since the state space contains every ghost number.
The part that survives is that the $r$-string pairing $\pair[r]{\cdot}{\cdot}$ has, in their grading, the length-dependent degree $- r - 2$. For $r = 1$, this is the value $- 3$ of the single-string pairing, which they describe as a form of degree three.
Therefore, none of these pairings can serve as a cyclic form on the whole tower. This does not exclude another form.

The cubic action evades the first obstruction at the classical level, by pairing each auxiliary field with a multiplier in the multiplier form, or by using the even forms $\vertex{2k}$, reduced, on the quotient by their kernels.
It does not supply an odd symplectic form of uniform degree on the multi-string space, which would present the tower as a cyclic $\Loo$ algebra on the fields alone. A formulation which requires such a form meets their first obstruction unchanged.

The second obstruction does not arise for the cubic action: there are no type-II fields, and the equations of motion are the gradients of an action.
The $\Psi$ equation contains $\vertex{4}(\Sigma_2, \Psi, \cdot)$ and the $\Sigma_2$ equation contains $\vertex{4}(\cdot, \Psi^2)$, whose derivatives with respect to $\Sigma_2$ and to $\Psi$ are transposes of each other.
We have checked in finite-dimensional models, of the kind described in \cref{p:app:finite-models}, that their canonical equations on $(\Psi, \Sigma_2)$ alone are the gradient of no action for any constant non-degenerate form.
An action is obtained either by enlarging the field space, as the multipliers do, or by changing the equations off the diagonal on the same fields, as the quartic vertex does in the cubic action.

\subsubsection{Limits of the comparison}
\label{p:sec:barcobar:limits}

The comparison is classical.
The formulation of the quantum theory on their enlarged space is left open in their paper, and we have not compared the loop expansions of the two constructions.
Their construction is algebraic and, under their assumptions, applies to any $\Loo$ field theory.
We carry out our transformation for the bosonic closed string, and we do not examine here whether the kernels of the vertices and the reality conditions carry over to the superstring.


\subsection{The two obstructions and the scalar example in detail}
\label{sec:bc:detail}

\status{mixed}\record{024-hs-vs-bar-cobar, sections 5 and 6 (R34)}

This section gives the analysis behind the paragraphs on the two obstructions in \cref{p:sec:barcobar:structure} and behind the scalar example.
Three restrictions hold throughout.
The finite-dimensional checks are made with commuting fields, so that every sign fixed below by the Grassmann parity is bookkeeping by hand.
The self-adjointness test concerns constant forms only: a form that depends on the fields, a curved symplectic structure on the type-I space, is not excluded, and they do not exclude it either.
The word \enquote{evades} refers to the cubic action and to the multiplier form as classical actions; where the first obstruction persists is said at the end of the section.

\subsubsection{First obstruction: the degree of the length-diagonal pairings}
\label{sec:bc:obstructions}

Their first obstruction is that the enlarged space admits no non-degenerate bilinear form of degree three, the degree that the cyclic form of an action $\pair{\Phi}{\mathcal E}$ must have, since the field $\Phi$ has degree $-1$ and the equation $\mathcal E = (B + \Delta_c) \Phi + \Phi \otimes' \Phi$ has degree $-2$.
With the homological sign of their degrees, which we keep, the form must pair two elements whose degrees add up to $-3$.
In their scalar example the argument is a counting one: every vector of the enlarged space has negative degree, so that nothing is available to pair against.
In closed string field theory the state space contains every ghost number, every degree is populated at every length, and the counting does not apply as such.

A sharper form of the obstruction survives.
By~\eqref{p:eq:barcobar:degree}, a type-I element of length $\ell$ and total ghost number $N$ has degree $2 \ell - 1 - N$.
The pairing that the theory supplies on the length-$\ell$ component is the $\ell$-string pairing $\pair[\ell]{\cdot}{\cdot}$, which is non-zero only when the ghost numbers of its two arguments add up to $5 \ell$.
Its degree is therefore the sum of the degrees of two elements of length $\ell$ whose ghost numbers add up to $5 \ell$:
\begin{equation}
  \label{eq:bc:pairing-degree}
  \deg \pair[\ell]{\cdot}{\cdot}
    = 2 \, (2 \ell - 1) - 5 \ell
    = - \ell - 2,
\end{equation}
which equals $-3$ only for $\ell = 1$.
Each length-diagonal pairing thus has the wrong degree at every length $\ell \ge 2$, and the mismatch $\ell - 1$ depends on the length, so that no single shift repairs all of them.
At $\ell = 2$ the mismatch is the one unit of the shifted two-string pairing of \cref{ch:conventions}, and for even $\ell$ it is also a mismatch of parity: the $\ell$-string pairing is Grassmann-even and graded-symmetric, while a form of odd degree must be odd.
This is the statement that the two-string pairing is graded-symmetric and not an odd symplectic form.
In our words, an action of the form $\pair{\Phi}{\mathcal E}$ on the type-I space would need a shift of $\pair[\ell]{\cdot}{\cdot}$ that depends on the length.
The count is a statement about these forms, not a proof that no other form exists.

\paragraph{How the multiplier form evades it.}
The multiplier form puts no form on the type-I space.
It pairs the length-$\ell$ equation of the chain, the constraint $\Sigma_\ell - \Sigma_{\ell - 1} \otimes \Psi$ of ghost number $2 \ell$, with a new field $\Theta_\ell$ in the dual space, through the $\ell$-string pairing used as a duality pairing, whatever its degree~\eqref{eq:bc:pairing-degree}.
Saturation of the pairing fixes $\Ngh(\Theta_\ell) = 5 \ell - 2 \ell = 3 \ell$ and the Grassmann parity $\ell$ of $\Theta_\ell$, so that each term $\pair[\ell]{\Theta_\ell}{\cdot}$ of the action is even, as an action term must be.
The multipliers are not elements of the enlarged space: their degree is $- \ell - 1$, not $-1$.
The cotangent construction asks nothing of the degree of the pairing, only that it be non-degenerate, and the $\ell$-string pairing is, because the single-string pairing is.
The only form of degree $-3$ that the multiplier form uses is the single-string pairing itself, on the one field that is not doubled, $\Psi$.
We do not claim a statement at the level of the Batalin--Vilkovisky formalism: the total pairing, the single-string pairing on $\Psi$ plus the $\ell$-string pairings between $\Sigma_\ell$ and $\Theta_\ell$, has a degree that depends on the sector, and whether the multiplier form is the classical part of a master action with one odd symplectic form of uniform degree is open.

\paragraph{How the cubic action evades it.}
The cubic action needs no odd pairing at all on the auxiliary sector, because $\Sigma_2$ is an auxiliary field without a $Q_B$ kinetic term.
Its quadratic term $- \tfrac{\gs^2}{4!} \, \vertex{4}(\Sigma_2^2)$ is a Grassmann-even, symmetric form supplied by the quartic vertex of the theory.
It saturates at ghost number $4 + 4 = 8$, so that it has degree $2 \cdot 3 - 8 = -2$ in their counting, and it is non-degenerate on the quotient of the $2$-string space by $\ker \sprod{2}{2}$.
The identity~\eqref{p:eq:barcobar:identity} is the bridge between the two presentations: the multiplier~\eqref{p:eq:barcobar:composite-multiplier} fixed by the $\Sigma_2$ equation converts the duality pairing of the multiplier form into the quadratic term of the cubic action.
At length $k$ the same holds with the quadratic form of $\Sigma_k$, built from $\vertex{2k}$; whether it is degenerate is discussed below.

\subsubsection{Second obstruction: self-adjointness across factor numbers}

Their second obstruction concerns a form diagonal in factor number and length.
Since $\Delta_c$ and $\otimes'$ raise the factor number by one, $\pair{\Phi_2}{\Phi_1 \otimes' \Phi_1}$ can be non-zero while the cyclically permuted $\pair{\Phi_1}{\Phi_2 \otimes' \Phi_1}$ vanishes, and $\pair{\Phi_f}{\Delta_c \Phi_{f-1}}$ can be non-zero while $\pair{\Delta_c \Phi_f}{\Phi_{f-1}}$ vanishes.
In other words, the equation at factor number $f$ involves the fields at factor numbers $f - 1$ and $f$ but not $f + 1$, so that the Jacobian of the equations is block-triangular in the factor number and cannot be self-adjoint for any form diagonal in it.

The precise criterion is the Helmholtz condition of the inverse problem of the calculus of variations.
A system of equations $\mathcal E(x) = 0$ is the set of Euler--Lagrange equations of an action, through a constant non-degenerate form $\Omega$, if and only if the matrix
\begin{equation}
  \label{eq:bc:helmholtz}
  \Omega \, \frac{\pd \mathcal E}{\pd x}
\end{equation}
is symmetric at every point $x$.
On the canonical sector, the condition fails already at the lowest level.
Take the pair formed by the one-output component of $B \Phi_1 = 0$ and the first factorization condition, on the fields $(\Psi, \Phi_{1, 2})$ with the canonical value $\Phi_{1, 2} = \Psi^2 / 2$:
\begin{equation}
  \mathcal E
    = \Big( Q_B \Psi + \gs \, \ell_2(\Phi_{1, 2}) + \cdots, \
      \Phi_{1, 2} - \tfrac{1}{2} \, \Psi^2 \Big).
\end{equation}
Its Jacobian is
\begin{equation}
  \label{eq:bc:jacobian}
  \frac{\pd \mathcal E}{\pd (\Psi, \Phi_{1, 2})}
    = \begin{pmatrix}
      Q_B & \gs \, \ell_2 \\
      - \Psi \wedge \cdot & 1
    \end{pmatrix},
\end{equation}
where $\Psi \wedge \cdot$ maps $\delta \Psi$ to the symmetrized product $\Psi \wedge \delta \Psi$.
Its two off-diagonal blocks are unrelated: $\gs \, \ell_2$ does not depend on $\Psi$ and $- \Psi \wedge \cdot$ does, so that no constant $\Omega$ makes~\eqref{eq:bc:helmholtz} symmetric.

We have tested the stronger form of this statement in a finite-dimensional model with commuting fields.
For the canonical system on $(\Psi, \Phi_{1, 2})$, the space of constant forms $\Omega$ for which~\eqref{eq:bc:helmholtz} is symmetric at three random points was computed by a singular value decomposition.
It was computed on the full tensor space $\Hs \oplus \Hs \otimes \Hs$, of dimension $d + d^2$ for $\dim \Hs = d$, rather than on $\Hs \oplus \Hn{2}$; the antisymmetric block decouples with a unit Jacobian and admits the identity form, so that the conclusion carries over to the symmetric space.
This space of forms is large, of dimension $29$ for $d = 3$ and $92$ for $d = 4$, because the $\Phi_{1, 2}$ block of the constraint is the identity and leaves the corresponding block of $\Omega$ nearly free.
But every element of it is degenerate: over random combinations the maximal rank is $10$ out of $12$ for $d = 3$ and $17$ out of $20$ for $d = 4$.
Their canonical equations are therefore the Euler--Lagrange equations of no constant non-degenerate form on the type-I fields, not only of no form diagonal in the factor number: the second obstruction is realized in the model in a sharpened form.
Nothing weaker than an enlargement of the field space evades it with a constant form.

\paragraph{How the multiplier form evades it.}
The multiplier form changes the system rather than finding a form for theirs.
The doubling adds the equations of the multipliers and makes every constraint the $\Theta_\ell$-gradient of a term $\pair[\ell]{\Theta_\ell}{\mathcal E_\ell}$, so that the doubled system is a gradient by construction.
Its Jacobian is self-adjoint for the cotangent form on $(\Psi; \Sigma_\ell, \Theta_\ell)$,
\begin{equation}
  \label{eq:bc:cotangent-form}
  \Omega_{\mathrm{cot}}
    = \pair{\cdot}{\cdot}
      \oplus \bigoplus_{\ell \ge 2}
      \begin{pmatrix}
        0 & \pair[\ell]{\cdot}{\cdot} \\
        \pair[\ell]{\cdot}{\cdot} & 0
      \end{pmatrix},
\end{equation}
which is off-diagonal between $\Sigma_\ell$ and $\Theta_\ell$: this is exactly what the block-triangular form requires, since the factorization equation is paired with a field outside the type-I space.
There are no type-II fields, so that no cyclicity across factor numbers is needed, and their product $\otimes'$ survives only as the bilinear term $\Sigma_{\ell - 1} \otimes \Psi$ of the constraint, symmetric by construction of $\Hn{\ell}$; this is why the vertices $\vertex{n}(\Sigma_n)$ are cyclic with no further condition.

We have checked this in the model, with the chain action~\eqref{p:eq:barcobar:chain-action} truncated at $k \le 2$ and with the cubic vertex.
Its Euler--Lagrange system, written by hand, has a Jacobian that is self-adjoint for~\eqref{eq:bc:cotangent-form}, to $10^{-8}$; the system written by hand is the gradient of the action, by an independent computation; and removing from the $\Psi$ equation the back-reaction of the multiplier, the term $\Theta_2 \cdot \Psi$, which is what their undoubled $\Psi$ equation looks like, destroys the symmetry, with an asymmetry of $3.6$.
Since the cyclicity of an action is automatic, this test is a check of the equations, not of the action: the back-reaction term is the transpose of the dependence of the constraint on $\Psi$.

\paragraph{How the cubic action evades it.}
The mechanism is the same with another form.
The pair $(\Psi, \Sigma_2)$ with the gradients of the cubic action is self-adjoint for the single-string pairing plus the two-string pairing, which we also checked in the model.
The block-triangular form is broken because the $\Psi$ equation now contains $\vertex{4}(\Sigma_2, \Psi, \cdot)$, which depends on $\Psi$, and the $\Sigma_2$ equation contains $\vertex{4}(\Sigma_2 - \Psi^2, \cdot)$: up to the factor fixed by~\eqref{p:eq:barcobar:identity}, both off-diagonal blocks are $\gs^2 \, \vertex{4}(\cdot, \Psi, \cdot)$.
The price, as in the truncation subsection, is that the $\Sigma_2$ equation imposes the chain only modulo the kernel of the quartic form.

\subsubsection{Where an obstruction persists}

The first obstruction does persist in one place, which is already part of the construction: the Batalin--Vilkovisky, odd-symplectic presentation of the type-I tower with the physical grading.
That the two-string pairing is graded-symmetric and even, so that the symplectic partner needs the shifted space $(\Hs^{\otimes 2})[1]$ (\cref{ch:conventions}), and that no pairing diagonal in the degree makes the products on the multi-string space cyclic (\cref{ch:ibl}), are by~\eqref{eq:bc:pairing-degree} the same fact as their first obstruction.
A Batalin--Vilkovisky action of the form $\pair{\Phi}{\mathcal E}$ on the type-I tower would need a shift $\ell - 1$ at length $\ell$, and no single form supplies it.
The cubic action is not affected, since it uses the even form $\vertex{4}$ and a classical action needs no odd form on its auxiliary fields; the multiplier form is not affected either, being a cotangent construction.
The statement that may be made is therefore that the construction supplies, at the classical level, the two things the bar--cobar equations lack for an action: a pairing, which is the quartic vertex itself or a dual tower of multipliers, and a field space on which the equations are self-adjoint, by keeping one cyclic physical field and no field of factor number two or more.
It does not provide a cyclic structure, an odd symplectic form of uniform degree, on the multi-string space.

The tower of all orders does not meet either obstruction.
Its Hessian is block-tridiagonal and symmetric because it is a Hessian: there are no type-II fields and no factor number to cross, and the off-diagonal blocks between $\Sigma_k$ and $\Sigma_{k+1}$ are transposes of each other by construction, the mechanism checked above for the multiplier form.
The first obstruction reappears only as the non-degeneracy of the quadratic form of each $\Sigma_k$ on the quotient of $\Hn{k}$ by the kernel of $\vertex{2k}$, which the continued-fraction elimination of \cref{ch:all-orders} needs, and which the quotient handles level by level.
In a model with random symmetric vertices there is no string kernel, so that a check of this point there tests nothing specific to the string; whether the dressing of order $\gs$ creates a kernel depending on $\Psi$ is a question about the poles of the elimination, not a structural obstruction of their kind.

\subsubsection{The scalar example}
\label{sec:bc:scalar}

Their example is the scalar theory
\begin{equation}
  \label{eq:bc:scalar-action}
  S
    = \int \dd x \left(
      \frac{1}{2} \, \phi \, (\pd^2 - m^2) \, \phi
      + \frac{g}{3} \, \phi^3
      + \frac{\lambda}{4} \, \phi^4
    \right),
  \qquad
  (\pd^2 - m^2) \, \phi + g \, \phi^2 + \lambda \, \phi^3
    = 0,
\end{equation}
with $\lambda$ the quartic coupling.
Its $A_\infty$ data have one degree-zero and one degree-$(-1)$ vector per point, the products $b_1 = \pd^2 - m^2$, $b_2 = g \, \delta$ and $b_3 = \lambda \, \delta \delta$, and no gauge symmetry, since there is no vector of degree one.
The type-I field is a sequence of functions $\phi_i(x_1, \ldots, x_i)$.
The equation $B \Phi = 0$ gives, for every $i$ and every slot $k \le i$,
\begin{equation}
  \label{eq:bc:scalar-B}
  (\pd_k^2 - m^2) \, \phi_i
    + g \, \phi_{i+1}(\ldots, x_k, x_k, \ldots)
    + \lambda \, \phi_{i+2}(\ldots, x_k, x_k, x_k, \ldots)
    = 0,
\end{equation}
and the factorization equation gives $\phi_{i+j}(x_1, \ldots, x_{i+j}) = \phi_i(x_1, \ldots, x_i) \, \phi_j(x_{i+1}, \ldots, x_{i+j})$ for all $i, j \ge 1$.
Its unique solution is $\phi_i = \phi_1(x_1) \cdots \phi_1(x_i)$, without the factor $1 / i!$ because their example uses the tensor normalization of the $A_\infty$ case.
On it every row of~\eqref{eq:bc:scalar-B} collapses to the equation of motion~\eqref{eq:bc:scalar-action} at the point $x_k$, multiplied by $\phi_1$ at the other points.

\paragraph{The multiplier form.}
With $\Sigma_k = \phi^{\otimes k}$ on the chain, matching their $\phi_i$, and the Euclidean pairing with one component per point, without ghost grading, the multiplier form of~\eqref{eq:bc:scalar-action} is
\begin{equation}
  \label{eq:bc:scalar-multiplier}
  S_{\mathrm{mult}}
    = \frac{1}{2} \int \phi \, (\pd^2 - m^2) \, \phi
      + \frac{g}{3} \int_x \Sigma_3(x, x, x)
      + \frac{\lambda}{4} \int_x \Sigma_4(x, x, x, x)
      + \sum_{k = 2}^4 \int \Theta_k \, (\Sigma_k - \Sigma_{k-1} \otimes \phi),
\end{equation}
with $\Sigma_1 = \phi$.
The vertices $\vertex{n}(\Sigma_n)$ are here the contact functionals $\int_x \Sigma_n(x, \ldots, x)$, the multilocal field evaluated on the diagonal: this is the scalar shadow of the multilocality of their type-I fields.
The Euler--Lagrange equations are solved from the top down.
The $\Theta_k$ equations impose the chain.
The $\Sigma_4$, $\Sigma_3$ and $\Sigma_2$ equations fix the multipliers:
\begin{equation}
  \Theta_4
    = - \frac{\lambda}{4} \, \delta_{x_1 x_2 x_3 x_4},
  \qquad
  \Theta_3
    = - \frac{g}{3} \, \delta_{x_1 x_2 x_3} + \Theta_4 \cdot \phi,
  \qquad
  \Theta_2
    = \Theta_3 \cdot \phi,
\end{equation}
where $\Theta_k \cdot \phi$ contracts the last slot of $\Theta_k$ with $\phi$.
The $\phi$ equation is
\begin{equation}
  \label{eq:bc:scalar-phi}
  (\pd^2 - m^2) \, \phi
    - 2 \, \Theta_2 \cdot \phi
    - \Theta_3(\Sigma_2, \cdot)
    - \Theta_4(\Sigma_3, \cdot)
    = 0
  \quad \Longrightarrow \quad
  (\pd^2 - m^2) \, \phi + g \, \phi^2 + \lambda \, \phi^3
    = 0.
\end{equation}
On the chain the three multiplier terms give $\tfrac{g}{3} \, (2 + 1) \, \phi^2 + \tfrac{\lambda}{4} \, (2 + 1 + 1) \, \phi^3$.
The $n$ slots of the contact vertex are recovered as $2$ from $\Theta_2$, whose constraint is quadratic in $\phi$ through $\Sigma_1 = \phi$, plus one for each higher multiplier, in total $n$, which cancels the $1/n$ in front of the vertex.
This is their collapse $\phi_i = \phi_1^{\otimes i}$ of~\eqref{eq:bc:scalar-B} reproduced from an action, with three multipliers, for their own example.

We have checked this on a periodic lattice of four sites in one dimension, with the lattice Laplacian for $\pd^2$.
The chain solves all the $\Sigma$ and $\Theta$ equations, to $10^{-10}$.
The $\phi$-gradient of~\eqref{eq:bc:scalar-multiplier}, computed by finite differences as an independent path, equals the collapsed equation of motion~\eqref{eq:bc:scalar-phi}, to $10^{-9}$.
Replacing $g / 3$ by $g / 4$ in $\Theta_3$ leaves a residual of $0.05$ in the $\Sigma_3$ equation, so that the check can fail.

\paragraph{The equation of motion term.}
The scalar theory has no gauge parameter, so that the term proportional to the equation of motion in their gauge transformation of $\Sigma_2$, \eqref{p:eq:barcobar:gauge-sigma}, has no scalar instance.
What survives is the structure behind it: the components of $B \Phi = 0$ with two or more output strands are proportional to the equation of motion on the chain, identities rather than equations.
For the scalar, the component $i = 2$, $k = 1$ reads, on the chain and for $\phi$ off-shell,
\begin{equation}
  \label{eq:bc:scalar-eom-row}
  (\pd_1^2 - m^2) \, \phi_2(x_1, x_2)
    + g \, \phi_3(x_1, x_1, x_2)
    + \lambda \, \phi_4(x_1, x_1, x_1, x_2)
    = \big[ (\pd^2 - m^2) \, \phi + g \, \phi^2 + \lambda \, \phi^3 \big](x_1) \, \phi(x_2),
\end{equation}
which we checked on the same lattice, to $10^{-15}$.
In the string case, their gauge variation of $\Sigma_2$ equals the composite law plus the equation of motion times $\Lambda$ for the same reason: $B$ acting on the chain produces the equation of motion in one slot times the rest of the chain, which is the factorized right-hand side of~\eqref{eq:bc:scalar-eom-row}, with the gauge parameter inserted in the remaining slot.
Without a gauge parameter, the factorization is all there is.

\paragraph{The cubic action.}
For contrast, the Hubbard--Stratonovich form of~\eqref{eq:bc:scalar-action} with a bilocal field $\Sigma_2(x, y)$ is
\begin{equation}
  S_{\mathrm{HS}}
    = \frac{1}{2} \int \phi \, (\pd^2 - m^2) \, \phi
      + \frac{g}{3} \int \phi^3
      - \frac{\lambda}{4} \int_x \Sigma_2(x, x)^2
      + \frac{\lambda}{2} \int_x \Sigma_2(x, x) \, \phi(x)^2.
\end{equation}
This is the quartic step of the construction with $\Hs$ replaced by functions on the lattice, and it was not checked by machine.
Its $\Sigma_2$ equation is $\Sigma_2(x, x) = \phi(x)^2$ on the diagonal only.
The quadratic form $\lambda \int_x \Sigma_2(x, x) \, \Sigma_2(x, x)$, which plays the role of $\vertex{4}$, is degenerate on every off-diagonal component of $\Sigma_2$: the scalar kernel of the quartic form is large, of codimension equal to the number of points in the space of symmetric bilocal fields, and the quotient by it reduces $\Sigma_2$ to the local field $\Sigma_2(x, x)$, the ordinary Hubbard--Stratonovich field.
In string field theory the quartic vertex is not a contact vertex, since it integrates over a region of the moduli space, and its kernel is small; this is why the string $\Sigma_2$ remains a bilocal field while the scalar one collapses to a local field.
